\chapter{Einleitung}

\section{Probabilistische Vorhersagen}
\raggedbottom
\glqq Kennt eine Intelligenz sämtliche Naturgesetze und den vollständigen Zustand der Welt, wäre sie dadurch in der Lage, Vergangenheit und Zukunft zu berechnen\grqq{}. Das ist die Idee des sogenannten \glqq Laplaceschen Dämons\grqq{} \parencite{deamon-out}, welcher auf den französischen Mathematiker Pierre-Simon Laplace zurückgeht \parencite{laplace1814essai}. Der Dämon nimmt dabei die Rolle einer allwissenden Entität ein. 

Hinter diesem Gedankenexperiment verbirgt sich neben einem streng deterministischen Weltbild auch ein grundlegendes menschliches Bedürfnis: zu wissen, was in der Zukunft geschehen wird. Dieses Bedürfnis offenbart sich konkret in zahlreichen Situationen. Unternehmen müssen abschätzen, wie sich Nachfrage und Preise entwickeln könnten, Banken berücksichtigen die erwartete Wirtschaftsentwicklung bei der Festlegung von Zinsen, Epidemiolog*innen beobachten mögliche Entwicklungen von Krankheitsgeschehen und in der Landwirtschaft entscheiden erwartete Veränderungen von Temperatur und Niederschlag über Anbau und Ernte. Fragen nach der Zukunft sind damit allgegenwärtig und für die Gesellschaft essenziell \parencite{prediction}.

Der Laplacesche Dämon lässt jedoch auf sich warten. Bereits die Annahme, dass die Welt vollständig deterministischen Gesetzmäßigkeiten folgt, ist nicht unumstritten \parencite{determinism}. Doch selbst unter der Annahme eines deterministischen Weltbildes folgt daraus nicht zwangsläufig eine vollständige Vorhersagbarkeit der Zukunft. Denn deterministische Systeme können aufgrund ihrer Dynamik schwer vorhersehbar sein \parencite{Shermer1995}. Insbesondere bei chaotischen Systemen wie dem Wetter können kleinste Unsicherheiten im Ausgangszustand zu erheblichen Unterschieden in der weiteren Entwicklung führen \parencite{dwd_ensemble_vorhersage}. So bleibt der Dämon letztlich eine Wunschvorstellung, eine Chimäre, wie \textcite{Shermer1995} festhält. Eine vollständige Kenntnis der Zukunft ist somit vorerst unerreichbar. Es bleibt uns also nichts anderes übrig, als auf Vorhersagen zurückzugreifen, die mit Unsicherheit behaftet sind. Vorhersagen stehen also in jenem Spannungsfeld zwischen dem menschlichen Wunsch nach Gewissheit und der tatsächlichen Unsicherheit der Zukunft.

Moderne mathematische Vorhersagen versuchen, zukünftige Entwicklungen auf Grundlage von Daten, Wahrscheinlichkeiten und überprüfbaren Modellen abzuschätzen \parencite{prediction}. Die zentrale Herausforderung besteht also nicht darin, die Zukunft mit absoluter Sicherheit vorherzusagen, sondern darin, die vorhandenen Informationen möglichst sinnvoll zu nutzen und die verbleibende Unsicherheit angemessen zu quantifizieren \parencite{wilks,gneiting2014probabilistic}. Dennoch wird diese Unsicherheit in der Praxis nicht immer explizit berücksichtigt. Stattdessen beschränken sich einige Vorhersagen auf punktuelle Aussagen oder vage Formulierungen, ohne Unsicherheiten systematisch zu quantifizieren \parencite{gneiting2014probabilistic}. Genau an dieser Stelle setzt diese Arbeit an.

Eine Möglichkeit, diese Unsicherheit explizit zu berücksichtigen, sind probabilistische Vorhersagen \parencite{gneiting2014probabilistic}. Dabei wird der Vorhersagegegenstand selbst als Zufallsvariable modelliert und eine Vorhersage durch eine (bedingte) Wahrscheinlichkeitsverteilung anstelle eines einzelnen Punktwertes beschrieben. Während eine deterministische Wettervorhersage beispielsweise lautet: \glqq Morgen wird es regnen\grqq{} oder \glqq Morgen werden es $15\C$\grqq{}, könnte eine probabilistische Vorhersage stattdessen lauten: \glqq Das Regenrisiko beträgt morgen $75\%$\grqq{} oder \glqq Die Temperatur morgen entspricht der Realisierung einer normalverteilten Zufallsvariablen mit Mittelwert $15\C$ und Standardabweichung $2\C$\grqq{}. Während die Angabe eines Regenrisikos im Alltag bereits üblich ist, erscheint die Beschreibung einer Temperatur durch eine Wahrscheinlichkeitsverteilung zunächst ungewohnt, obwohl beide Vorhersagen auf demselben probabilistischen Prinzip beruhen.

Auf rein deterministische Aussagen reduziert, können verschiedene probabilistische Vorhersagen identisch erscheinen. So implizieren sowohl ein vorhergesagtes Regenrisiko von $51\%$ als auch von $90\%$, dass bei einer Entscheidungsregel mit einem Schwellenwert von $50\%$ die deterministische Aussage \glqq Morgen wird es regnen\grqq{} getroffen wird. Ebenso führen zwei Temperaturvorhersagen mit demselben Mittelwert von $15\C$, aber unterschiedlichen Standardabweichungen zur gleichen Punktvorhersage von $15\C$. Der Unterschied liegt jedoch in der quantifizierten Unsicherheit. Eine größere Standardabweichung bedeutet, dass ein breiterer Bereich möglicher Temperaturen plausibel ist. Diese zusätzliche Information ermöglicht es, Entscheidungen an die Unsicherheit der Vorhersage anzupassen, etwa vorsorglich einen zusätzlichen Pullover einzupacken oder Aktivitäten flexibler zu planen.

Durch die explizite Quantifizierung von Unsicherheit können probabilistische Vorhersagen daher prinzipiell bessere Entscheidungen ermöglichen. Die Quantifizierung von Unsicherheit bedeutet jedoch nicht, dass eine probabilistische Vorhersage auch eine gute Vorhersage ist. Auch probabilistische Vorhersagen können systematische Fehler aufweisen oder ihre Unsicherheit unangemessen abbilden. Für den praktischen Nutzen einer Vorhersage ist es daher ebenso wichtig zu beurteilen, wie gut sie tatsächlich ist und welche Fehler und Schwachstellen mit ihr verbunden sind.

Diese systematische Beurteilung der Güte von Vorhersagen ist Gegenstand der Vorhersageevaluation oder Verifikation \parencite{wilks}. Für probabilistische Vorhersagen stellt dies eine besondere Herausforderung dar, weil nicht einzelne vorhergesagte Werte, sondern vollständige Vorhersageverteilungen hinsichtlich ihrer Güte bewertet werden müssen. Zudem liegen in der Regel nur endlich viele Beobachtungen eines möglicherweise zeitlich nicht stationären Systems vor. Dabei ist insbesondere von Interesse, ob die vorhergesagten Wahrscheinlichkeiten mit der tatsächlichen Häufigkeit des Eintretens übereinstimmen und ob die Unsicherheit der Vorhersagen angemessen abgebildet wird.

Die Evaluation der Güte probabilistischer Vorhersagen ist daher nicht trivial. Sie stellt das zentrale Thema der vorliegenden Arbeit dar. Dabei betrachten wir das Thema im Hinblick auf Temperaturvorhersagen, weswegen wir uns auf die Evaluation kontinuierlicher Vorhersagegegenstände fokussieren. Dazu werden in der vorliegenden Arbeit sowohl aktuelle als auch klassische Evaluationsmethoden probabilistischer Vorhersagen theoretisch vorgestellt, implementiert und angewendet. Ziel ist es, einen Überblick über verschiedene Möglichkeiten zur Beurteilung probabilistischer Vorhersagen zu geben und damit eine Orientierung dafür zu schaffen, wie eine Evaluation aufgebaut werden kann, welche Methoden sich für die Untersuchung unterschiedlicher Aspekte der Vorhersagegüte eignen und wie diese angewendet werden können. Die Ziele der Arbeit fassen wir im nächsten Abschnitt zusammen.

\section{Ziele der Arbeit}
Diese Arbeit gliedert sich in einen theoretischen und einen praktischen Teil. Der theoretische Teil behandelt stochastische Grundlagen und ausgewählte statistische Methoden zur Evaluation probabilistischer Vorhersagen für kontinuierliche Vorhersagegegenstände. Der Inhalt orientiert sich dabei maßgeblich an \textcite{gneiting2023regression} sowie \textcite{gneiting2014probabilistic}, wobei auf dieser Grundlage eigene Schwerpunkte gesetzt und ausgewählte Aspekte durch eigene Illustrationen vertieft werden. Ziel ist es, einen mathematischen Überblick über die Evaluation probabilistischer Vorhersagen zu geben und die theoretischen Grundlagen der verwendeten Evaluationsmethoden zu vermitteln.  

Dazu werden in Kapitel \ref{chap:Grundlagen} zunächst die mathematische Struktur sowie die grundlegenden Begriffe zur Evaluation probabilistischer Vorhersagen eingeführt. In den Kapiteln \ref{chap:classic_calibration}, \ref{chap:quant_cal} und \ref{chap:cal_relation} widmen wir uns anschließend dem Aspekt der Kalibrierung, wobei wir neben klassischen Kalibrierungsbegriffen auch eine auf Quantilen basierende Kalibrierung betrachten. Der zweite Schwerpunkt der Arbeit liegt auf der Bewertung probabilistischer Vorhersagen mithilfe von Scores, die in Kapitel \ref{chap:scoring} methodisch eingeführt und hinsichtlich ihrer Eigenschaften betrachtet werden. Abschließend werden in Kapitel \ref{chap:score_decomp} Möglichkeiten zur Zerlegung dieser Scores vorgestellt, die eine differenzierte Betrachtung der Güte probabilistischer Vorhersagen ermöglichen. 

Im praktischen Teil der Arbeit werden die im theoretischen Teil vorgestellten Methoden zur Evaluation und zum Vergleich probabilistischer Vorhersagen auf ein konkretes Beispiel angewendet. Zunächst beschreiben wir dazu in Kapitel \ref{chap:nwp} allgemein numerische Wettervorhersagemodelle. Anschließend werden in Kapitel \ref{chap:data} die für die vorliegende Arbeit ausgewählten Daten beschrieben. Diese umfassen die Temperaturvorhersagen des numerischen Wettervorhersagemodells ICON-D2-EPS \parencite{reinert2025icon_dwd} sowie die Beobachtungsdaten des HOSTRADA \parencite{krahenmann_high-resolution_2018}. Dabei betrachten wir die 24- und 48-Stunden Vorhersagen für den Zeitraum November 2025. In Kapitel \ref{chap:meth} beschreiben wir das Evaluationsvorgehen. Dabei erläutern wir insbesondere, wie aus den ICON-D2-EPS-Modellvorhersagen probabilistische Vorhersagen erstellt werden können. Anschließend untersuchen wir die verschiedenen Vorhersagen anhand der im theoretischen Teil eingeführten Kalibrierungs- und Bewertungsmaße für die voliegenden Daten. Die Ergebnisse dieser Evaluation werden in Kapitel \ref{chap:res} dargestellt. Abschließend diskutieren wir in Kapitel \ref{chap:discuss} die Ergebnisse des praktischen Teils und reflektieren den Nutzen der verwendeten Evaluationsmethoden.

\section{Hinweis zum bereitgestellten Material}
\label{sec:material}
Der vollständige Code, die Quelldateien dieser Arbeit sowie die für die Reproduktion der Analysen erforderlichen Dateien sind über das GitHub-Repository
\begin{quote}
    \url{https://github.com/harro-rejewski/bachelor_thesis.git}
\end{quote}
verfügbar. Weitere Informationen zur Implementierung finden sich in Abschnitt \ref{sec:sofware}.
